% GENERATED FILE. Edit canonical lessons, then run npm run generate:books.
% canonical-source-hash: fnv1a64:29ea61d4dfc74c11
% canonical-lessons: ML-C11-nirangal, ML-S114-letter-ma

\chapter{Four Core Colours}
\label{ch:ml-core-colours}
\hlchaptermodality{11}

\begin{chapteropening}
\textbf{By the end of this chapter:} \emph{I can name black, white, red and blue in Malayalam.}

Say karuppŭ, veḷḷa, cuvappŭ and nīla, and mark which three are native Dravidian and which one is a shared Sanskrit loan.
\end{chapteropening}

\section[\texorpdfstring{\ml{കറുപ്പ്} \ml{വെള്ള} \ml{ചുവപ്പ്} \ml{നീല} (karuppŭ veḷḷa …)}{കറുപ്പ് വെള്ള ചുവപ്പ് നീല (karuppŭ veḷḷa …)}]{\ml{കറുപ്പ്}, \ml{വെള്ള}, \ml{ചുവപ്പ്}, \ml{നീല} — the same split as Tamil}

\label{lesson:ML-C11-nirangal}



\begin{quote}
You've seen Malayalam and Tamil share day-names almost word-for-word. Their colors follow the exact same family pattern.
\end{quote}

\subsection*{What to know first}
\begin{itemize}
\raggedright
  \item \textbf{\ml{കറുപ്പ്}} (\emph{karuppŭ}) = \textquotedblleft{}\textbf{black}\textquotedblright{} — matches Tamil \emph{karuppu} almost exactly.
  \item \textbf{\ml{വെള്ള}} (\emph{veḷḷa}) = \textquotedblleft{}\textbf{white}\textquotedblright{} — the same \textbf{veḷ-} \textquotedblleft{}bright/silvery\textquotedblright{} root as Tamil \emph{veḷḷai}.
  \item \textbf{\ml{ചുവപ്പ്}} (\emph{cuvappŭ}) = \textquotedblleft{}\textbf{red}\textquotedblright{} — a related word from the same Dravidian root family as Tamil \emph{sivappu}, though the two differ in both their opening consonant's sound and their vowel — a looser cousin than \emph{karuppŭ}/\emph{karuppu} or \emph{veḷḷa}/\emph{veḷḷai}, which match almost exactly.
\end{itemize}

\begin{cousinweb}
\textbf{\ml{നീല}} (\emph{nīla}) = \textquotedblleft{}\textbf{blue}\textquotedblright{} — the same \textbf{Sanskrit} loan Tamil uses (\emph{nīlam}), and the same one Hindi uses (\emph{nīlā}). Blue is where Malayalam, Tamil, and Hindi all converge on one shared Sanskrit word, even though Malayalam and Tamil otherwise built their own native words for black, white, and red.
\end{cousinweb}

\subsection*{Your turn}
Say these aloud:

\begin{itemize}
\raggedright
  \item \textquotedblleft{}karuppŭ, veḷḷa, cuvappŭ\textquotedblright{} — black, white, red
  \item \textquotedblleft{}nīla\textquotedblright{} — blue, shared with Tamil and Hindi
  \item the family match — \textquotedblleft{}karuppŭ/karuppu, veḷḷa/veḷḷai, cuvappŭ/sivappu\textquotedblright{}
\end{itemize}

\begin{tcolorbox}[breakable,colback=teal!4,colframe=teal!35!black,title={Before you move on}]
Which three Malayalam colors closely match Tamil's own words? (\textbf{Karuppŭ, veḷḷa, cuvappŭ} — black, white, red.) Which color is a shared Sanskrit loan across Malayalam, Tamil, and Hindi? (\textbf{Blue} — \ta{நீல} \emph{nīla} / \ta{நீலம்} \emph{nīlam} / \dv{नीला} \emph{nīlā}.)
\end{tcolorbox}
\section[\texorpdfstring{\ml{മ} (ma)}{മ (ma)}]{\ml{മ} — one character, met inside words you already say}

\label{lesson:ML-S114-letter-ma}



\begin{quote}
Before the new one: \ml{ന} — what does it do?

One character this time. Just one — and you have been saying it for pages without knowing which mark on the page it was.
\end{quote}

\subsection*{Script you'll notice: \ml{മ}}
\textbf{\ml{മ}} — \emph{ma}.

It is a \textbf{consonant}, and in this script a consonant is never bare: it comes with an \emph{a} already in it. So it is not \emph{m}, it is \textbf{ma}.

You already say these, and every one of them has \ml{മ} somewhere inside it:

\begin{itemize}
\raggedright
  \item \textbf{\ml{നമസ്കാരം}} \emph{namaskāraṁ} — hello / greetings (namaskāraṁ — \textquotedblleft{}a making of a bow\textquotedblright{})
  \item \textbf{\ml{സാരമില്ല}} \emph{sāramilla} — it doesn't matter / no problem / you're welcome
  \item \textbf{\ml{താമസിക്കുക}} \emph{tāmasikkuka} — to live, to stay
  \item \textbf{\ml{എനിക്ക്} \ml{മലയാളം} \ml{അറിയാം}} \emph{enikkŭ malayāḷaṁ aṟiyāṁ} — 'I know Malayalam'
\end{itemize}

\subsection*{Writing: \ml{മ} — follow the numbered strip}
\hlblockfigure{\detokenize{figures/ML-S114-letter-ma-filmstrip.pdf}}{How \ml{മ} is written, stroke by stroke}

Follow the numbered strip: start where movement 1 begins, and let each panel show you the next movement before your pen makes it. Then write \ml{മ} yourself — slowly, and larger than it is printed.

\begin{quote}
The strip shows \textbf{where to start the character and which way to travel}, in one attested order. That is taught with real variation from school to school, so the strip names its source beneath it: treat it as a sound way in, not the only one.
\end{quote}

\subsection*{Your turn}
\begin{itemize}
\raggedright
  \item \emph{Look:} at these words, and find \ml{മ} in the ones that have it
\end{itemize}

\begin{quote}
\ml{നമസ്കാരം}  ·  \ml{സാരമില്ല}  ·  \ml{നന്ദി}
\end{quote}

\begin{itemize}
\raggedright
  \item \emph{Trace it:} \ml{മ} three times, saying \emph{ma} as you finish each one
  \item \emph{Look:} back at any page of this chapter and find \ml{മ} once more
\end{itemize}

\begin{tcolorbox}[breakable,colback=teal!4,colframe=teal!35!black,title={Before you move on}]
Which character is this — \ml{മ}? What sound does it carry? (\emph{\textbf{ma}}.) Name one word you already say that contains it.
\end{tcolorbox}
